% ============================================================
% Section 51: 布里渊散射——光子与声子的非弹性散射
% ============================================================
\section{固体理论：布里渊散射——光子与声子的非弹性散射}
\label{sec:brillouin-scattering}

\begin{modelbox}[题目：光受透明固体的散射与声子相互作用]
光受透明固体的散射可以看成是光量子与声子相互作用过程的结果。试利用能量和动量守恒定律证明，由于散射而偏转了 $\theta$ 角的光束必定包含这样的成分，其波长相对于入射光移动了
\begin{equation}
\Delta\lambda = \pm\frac{2nv}{c}\lambda_0\sin\frac{\theta}{2}
\end{equation}
式中，$n$ 为物质的折射率，$v$ 为固体中的声速，$c$ 为真空中的光速，$\lambda_0$ 为入射光在真空中的波长。
\end{modelbox}

\begin{tipbox}[考点快速分析]
\begin{itemize}
    \item \textbf{布里渊散射}：光子与声子的非弹性散射，属于拉曼散射的近亲，频移在 $10^{9}\sim10^{11}\,\mathrm{Hz}$ 量级
    \item \textbf{守恒律}：准动量守恒 $\vec k_0 = \vec k \pm \vec q$（$+$ 发射声子，$-$ 吸收声子）；能量守恒 $\hbar\omega_0 = \hbar\omega \pm \hbar\omega_\mathrm{ph}$
    \item \textbf{小频移近似}：声子能量远小于光子能量，$|\omega_0 - \omega| \ll \omega_0$，故 $|\vec k| \approx |\vec k_0|$
    \item \textbf{斯托克斯 / 反斯托克斯}：发射声子（$+$）对应斯托克斯线（红移）；吸收声子（$-$）对应反斯托克斯线（蓝移）
\end{itemize}
\end{tipbox}

\begin{derivationbox}[标准解题过程]
\textbf{Step 1: 光子与声子的动量关系}

光在介质中的波矢大小为
\begin{equation}
|\vec k| = \frac{2\pi n}{\lambda},
\end{equation}
其中 $\lambda$ 为光在真空中的波长。对于可见光，光子动量远大于声子动量，散射过程中光子能量变化极小，因此可近似认为散射光子动量大小等于入射光子动量：
\begin{equation}
|\vec k| \approx |\vec k_0| = \frac{2\pi n}{\lambda_0}.
\end{equation}

设入射光波矢为 $\vec k_0$，散射光波矢为 $\vec k$，散射角为 $\theta$。在声子参与的过程中，准动量守恒为
\begin{equation}
\vec k_0 = \vec k \pm \vec q,
\end{equation}
其中 $\vec q$ 为被吸收（$-$）或发射（$+$）的声子波矢。由矢量几何关系（等腰三角形，顶角为 $\theta$），声子波矢大小为
\begin{equation}
|\vec q| = 2|\vec k_0|\sin\frac{\theta}{2} = \frac{4\pi n}{\lambda_0}\sin\frac{\theta}{2}.
\label{eq:brillouin-q}
\end{equation}

\textbf{Step 2: 能量守恒与频率关系}

设声速为 $v$，声子角频率为 $\omega_\mathrm{ph} = v|\vec q|$（声学支色散）。代入 \eqref{eq:brillouin-q}：
\begin{equation}
\omega_\mathrm{ph} = v\cdot\frac{4\pi n}{\lambda_0}\sin\frac{\theta}{2} = \frac{4\pi nv}{\lambda_0}\sin\frac{\theta}{2}.
\end{equation}

入射光子和散射光子的角频率分别为
\begin{equation}
\omega_0 = \frac{2\pi c}{\lambda_0}, \qquad \omega = \frac{2\pi c}{\lambda}.
\end{equation}
（注：光子的能量仅与频率有关，与介质折射率无关，故用真空中的 $c$ 和真空波长。介质中的波矢因子 $n$ 已在动量守恒中体现。）

能量守恒定律为
\begin{equation}
\hbar\omega - \hbar\omega_0 = \pm\hbar\omega_\mathrm{ph},
\end{equation}
正号对应\textbf{反斯托克斯过程}（吸收声子，散射光频率升高），负号对应\textbf{斯托克斯过程}（发射声子，散射光频率降低）。代入频率表达式：
\begin{equation}
\frac{2\pi c}{\lambda} - \frac{2\pi c}{\lambda_0} = \pm\frac{4\pi nv}{\lambda_0}\sin\frac{\theta}{2}.
\end{equation}

两边约去 $2\pi c$：
\begin{equation}
\frac{1}{\lambda} - \frac{1}{\lambda_0} = \pm\frac{2nv}{c\lambda_0}\sin\frac{\theta}{2}.
\end{equation}

\textbf{Step 3: 波长移动的化简}

将左边通分：
\begin{equation}
\frac{\lambda_0 - \lambda}{\lambda\lambda_0} = \pm\frac{2nv}{c\lambda_0}\sin\frac{\theta}{2}.
\end{equation}

由于声子能量远小于光子能量，波长变化极小，$\lambda \approx \lambda_0$，于是 $\lambda\lambda_0 \approx \lambda_0^2$：
\begin{equation}
\frac{\lambda_0 - \lambda}{\lambda_0^2} \approx \pm\frac{2nv}{c\lambda_0}\sin\frac{\theta}{2}
\quad\Longrightarrow\quad
\Delta\lambda = \lambda - \lambda_0 \approx \pm\frac{2nv}{c}\lambda_0\sin\frac{\theta}{2}.
\end{equation}

证毕。
\end{derivationbox}

\begin{formulabox}[答案汇总]
\begin{center}
\begin{tabular}{cl}
\toprule
量 & 结果 \\
\midrule
声子波矢大小 $|\vec q|$ & $\displaystyle \frac{4\pi n}{\lambda_0}\sin\frac{\theta}{2}$ \\[10pt]
声子角频率 $\omega_\mathrm{ph}$ & $\displaystyle \frac{4\pi nv}{\lambda_0}\sin\frac{\theta}{2}$ \\[10pt]
波长移动量 $\Delta\lambda$ & $\displaystyle \pm\frac{2nv}{c}\lambda_0\sin\frac{\theta}{2}$ \\[10pt]
\bottomrule
\end{tabular}
\end{center}
\end{formulabox}

\begin{insightbox}[物理图像与概念深化]
\textbf{1. 布里渊散射与瑞利散射、拉曼散射的区别}

\begin{center}
\begin{tabular}{lccc}
\toprule
特征 & \textbf{瑞利散射} & \textbf{布里渊散射} & \textbf{拉曼散射} \\
\midrule
散射中心 & 静态杂质/密度涨落 & \textbf{声子}（声学支） & 光学声子 / 分子振动 \\
频移大小 & $0$（弹性） & $\sim 10^{9}\text{--}10^{11}\,\mathrm{Hz}$ & $\sim 10^{12}\text{--}10^{13}\,\mathrm{Hz}$ \\
频移来源 & --- & 声子能量 $\hbar\omega_\mathrm{ph}$ & 光学模 / 分子振动能级 \\
\bottomrule
\end{tabular}
\end{center}

布里渊散射本质上是光子与\textbf{声学声子}的非弹性碰撞。声子作为晶格集体运动的量子化激发，为光子提供了一个``运动靶''，使得散射光子的能量和动量发生可测量的改变。

\textbf{2. 斯托克斯线与反斯托克斯线的强度差异——为什么一个是 $n+1$、一个是 $n$？}

\textit{核心图像}：把晶格的每个振动模式看作一个量子谐振子，能级为 $E_n=(n+\tfrac{1}{2})\hbar\omega$。光子与谐振子相互作用时，通过升降算符改变声子数：

\begin{itemize}
    \item \textbf{发射声子}（斯托克斯）：谐振子从 $|n\rangle\to|n+1\rangle$，需要升算符 $a^\dagger$。由于 $a^\dagger|n\rangle=\sqrt{n+1}\,|n+1\rangle$，跃迁矩阵元 $|\langle n+1|a^\dagger|n\rangle|^2=n+1$。\textbf{即使 $n=0$（绝对零度），$+1$ 项仍然存在}——这叫``自发''发射，来源于量子力学的零点涨落。
    \item \textbf{吸收声子}（反斯托克斯）：谐振子从 $|n\rangle\to|n-1\rangle$，需要降算符 $a$。由于 $a|n\rangle=\sqrt{n}\,|n-1\rangle$，跃迁矩阵元 $|\langle n-1|a|n\rangle|^2=n$。\textbf{若 $n=0$，矩阵元为零}——没有声子可供吸收，过程无法进行。
\end{itemize}

在热平衡下，声子占据数服从玻色-爱因斯坦分布：
\begin{equation}
n_\mathrm{BE}(\omega_\mathrm{ph}) = \frac{1}{e^{\hbar\omega_\mathrm{ph}/k_\mathrm{B}T} - 1}.
\end{equation}

因此：
\begin{itemize}
    \item 斯托克斯强度 $\propto n_\mathrm{BE}+1$（包含自发 + 受激两部分）
    \item 反斯托克斯强度 $\propto n_\mathrm{BE}$（只有受激吸收）
\end{itemize}

\begin{center}
\begin{tabular}{ccc}
\toprule
温度 & 斯托克斯线 & 反斯托克斯线 \\
\midrule
$T=0$ & 存在（$n=0$，自发散射） & \textbf{消失}（$n=0$，无法吸收） \\
室温 & 强 & 较弱 \\
高温 & 都强，但斯托克斯始终更强 & 强度比 $I_\text{AS}/I_\text{S}=n/(n+1)\to1$ \\
\bottomrule
\end{tabular}
\end{center}

测量斯托克斯与反斯托克斯的强度比，可直接反解 $n_\mathrm{BE}$，从而测定样品温度——这是非接触式测温的重要方法。

\textbf{3. 红移与蓝移的物理——从光谱图理解}

``红移''和``蓝移''是光谱学里的标准说法，来源于光在光谱上的移动方向：

\begin{center}
\begin{tabular}{cccccc}
\toprule
紫光 & 蓝光 & 绿光 & 黄光 & 橙光 & \textbf{红光} \\
$\lambda$ 短 & $\to$ & $\to$ & $\to$ & $\to$ & $\lambda$ \textbf{长} \\
$\omega$ 高 & $\to$ & $\to$ & $\to$ & $\to$ & $\omega$ \textbf{低} \\
\bottomrule
\end{tabular}
\end{center}

\begin{itemize}
    \item \textbf{斯托克斯}（发射声子）：光子损失能量，$\omega$ 降低 → 光谱线向\textbf{红色端}移动 → \textbf{红移}
    \item \textbf{反斯托克斯}（吸收声子）：光子获得能量，$\omega$ 升高 → 光谱线向\textbf{蓝/紫端}移动 → \textbf{蓝移}
\end{itemize}

\textit{判断方法}：直接从能量守恒看符号。斯托克斯过程中光子能量减少：
\begin{equation}
\hbar\omega = \hbar\omega_0 - \hbar\omega_\mathrm{ph} \quad\Longrightarrow\quad \omega < \omega_0 \quad\Longrightarrow\quad \lambda > \lambda_0,
\end{equation}
波长变长，光谱向红端移动。反斯托克斯则相反。

\textbf{4. 为什么 $|\vec k|\approx|\vec k_0|$？}

因为声速 $v$ 远小于光速 $c$（典型值 $v\sim10^3\,\mathrm{m/s}$，$c\sim10^8\,\mathrm{m/s}$），声子能量 $\hbar\omega_\mathrm{ph}$ 只是光子能量 $\hbar\omega_0$ 的极小一部分：
\begin{equation}
\frac{\Delta\omega}{\omega_0}=\frac{\omega_\mathrm{ph}}{\omega_0}\sim\frac{v}{c}\sim10^{-5}.
\end{equation}
光子只``施舍''给声子约 $10^{-5}$ 的能量，自身波矢大小几乎不变。这正是弹性近似成立的基础。

\textbf{5. 布里渊散射的实验应用：测量声速}

从布里渊频移公式反解声速：
\begin{equation}
v = \frac{c\Delta\lambda}{2n\lambda_0\sin(\theta/2)} = \frac{c\Delta f}{2nf_0\sin(\theta/2)}.
\end{equation}

通过精确测量散射光的频率移动 $\Delta f$，可以无损伤地测定透明固体中的声速，进而研究其\textbf{弹性模量}、\textbf{相变}（声速在相变点附近反常）、以及\textbf{薄膜}和\textbf{微结构}的力学性质。布里渊光谱仪是现代凝聚态物理和材料科学的重要工具。

\textbf{6. 与 X 射线/中子非弹性散射的对照}

\begin{itemize}
    \item \textbf{X 射线非弹性散射}：光子能量高（$\sim 10^{4}\,\mathrm{eV}$），声子能量相对可忽略，但动量匹配好，可测完整声子色散 $\omega(q)$
    \item \textbf{中子非弹性散射}：中子质量大，$E\sim\hbar^2k^2/2M$，能量和动量都与声子匹配，是测量声子色散的``黄金标准''
    \item \textbf{布里渊散射}（可见光）：光子能量远大于声子，只能测 $q\to0$ 的长波声学支（$\omega=vq$），但设备简单、无损伤、空间分辨率高
\end{itemize}

三者的共同点是都利用\textbf{能量-动量守恒}来探测声子激发，区别在于探针与声子的能量/动量匹配程度不同。
\end{insightbox}

\begin{thinkbox}[考场速记：布里渊散射问题的四步法]
遇到``光与声子散射导致波长/频率移动''类问题：
\begin{enumerate}
    \item \textbf{写准动量守恒}：$\vec k_0 = \vec k \pm \vec q$，利用 $|\vec k|\approx|\vec k_0|$ 和散射角 $\theta$ 求 $|\vec q| = 2k_0\sin(\theta/2)$
    \item \textbf{写能量守恒}：$\hbar\omega_0 = \hbar\omega \pm \hbar\omega_\mathrm{ph}$，代入声子色散 $\omega_\mathrm{ph}=vq$
    \item \textbf{化简}：利用 $\lambda\approx\lambda_0$（小频移近似），通分后取 $\lambda\lambda_0\to\lambda_0^2$
    \item \textbf{定号}：斯托克斯（$+$，红移）vs 反斯托克斯（$-$，蓝移），注意温度对强度的影响
\end{enumerate}
\textit{记忆口诀}：动量守恒得波矢，能量守恒得频移，小量近似化波长，斯托克斯反斯托克斯分左右。
\end{thinkbox}
